problem:  
Let \((\mathbb{H}^n,\Delta,g_0)\) be the standard Heisenberg group with Carnot–Carathéodory distance \(d_{cc}\) and Haar measure \(\mu\). It is known that \((\mathbb{H}^n,d_{cc},\mu)\) satisfies the *sharp* measure contraction property \(\mathrm{MCP}(0,2n+3)\)—the best possible synthetic dimension among all \(N\).  
For a smooth function \(\psi:\mathbb{H}^n\to\mathbb{R}\) bounded from above and below, define the weighted measure \(\mu_\psi = e^{-\psi}\,\mu\).  

**Question:**  
Is the Haar measure \(\mu\) the **unique smooth measure (up to constant multiples)** on \(\mathbb{H}^n\) such that \((\mathbb{H}^n,d_{cc},\mu_\psi)\) satisfies the *optimal* measure contraction property \(\mathrm{MCP}(0,2n+3)\)?  
Equivalently, can any nonconstant smooth weight \(\psi\) preserve the sharp volume contraction inequality along sub-Riemannian geodesics, or does any nontrivial density necessarily increase the synthetic dimension \(N_{\mathrm{eff}}>2n+3\)?  

If density perturbations do spoil the optimal exponent, can one express the infinitesimal failure of MCP(0,2n+3) in terms of the horizontal gradient \(\nabla_H\psi\) and the geodesic Jacobian \(J\) of the sub-Riemannian exponential map (for instance, \( \frac{d}{ds}\log J + \nabla_H\psi\cdot \dot{\gamma}_H  \))?  

Why is it a "good" problem:  
This problem pinpoints a fundamental rigidity question at the intersection of sub-Riemannian geometry and metric measure theory: is the **Haar measure uniquely determined** by the requirement that it realizes the sharp synthetic dimensional bound of \(\mathrm{MCP}(0,2n+3)\)?  

The problem is compelling for several reasons:  
- **Conceptual simplicity:** the statement is purely geometric and measure-theoretic—yet it directly probes the deepest structure of the Heisenberg group’s synthetic curvature.  
- **Structural depth:** the sharp MCP exponent arises from the precise Jacobian of the Heisenberg exponential map. Determining whether any smooth weight preserves this inequality tests the fine geometry of horizontal geodesic flows and the group’s symmetries.  
- **Bridge across fields:** resolution would link measure contraction theory (optimal transport and synthetic curvature) with sub-Riemannian analysis (the geometry of nilpotent Lie groups, horizontal gradients, exponential maps). It would also touch weighted geometric inequalities reminiscent of Bakry–Émery curvature on singular spaces.  
- **Potential impact:** an affirmative rigidity result would give a deep characterization of the Haar measure as the *unique* measure achieving synthetic nonnegative curvature with minimal dimension—establishing a new kind of “synthetic uniqueness theorem.” Conversely, a counterexample would reveal unknown flexibility of sub-Riemannian volume geometry under weighting.  

Thus, though elementary to state, this problem sits precisely where analytic, geometric, and synthetic curvature theories converge, exemplifying the clarity, depth, and interdisciplinarity of a genuinely *good* research problem.